\section{动力学模型与问题描述}
\label{sec:model}

\subsection{等效关节空间模型}
\label{sec:modeling}

设基座体坐标系中的广义速度 \(\dot x_b=[v_b^{\top},\omega_b^{\top}]^{\top}\in\R^{6}\)（此处 \(\dot x_b\) 仅记速度，不表示姿态坐标的普通导数），关节变量 \(q\in\R^{n}\)。以下惯性块与雅可比均在一致的基座体坐标系表达，并可依赖基座姿态 \(g_b\)；在该表达下若惯性块与姿态无关，才简记为 \(H(q)\)。系统动能为二次型
\begin{equation}
T=\tfrac12\begin{bmatrix}\dot x_b^{\top}&\dot q^{\top}\end{bmatrix}
\begin{bmatrix}H_b&H_{bm}\\ H_{bm}^{\top}&H_m\end{bmatrix}
\begin{bmatrix}\dot x_b\\ \dot q\end{bmatrix},
\label{eq:kinetic}
\end{equation}
其中 \(H_b\in\R^{6\times6}\)、\(H_m\in\R^{n\times n}\)、\(H_{bm}\in\R^{6\times n}\)。由拉格朗日方程得完整动力学
\begin{equation}
\begin{bmatrix}H_b&H_{bm}\\ H_{bm}^{\top}&H_m\end{bmatrix}
\begin{bmatrix}\ddot x_b\\ \ddot q\end{bmatrix}
+\begin{bmatrix}C_b&C_{bm}\\ C_{bm}^{\top}&C_m\end{bmatrix}
\begin{bmatrix}\dot x_b\\ \dot q\end{bmatrix}
=\begin{bmatrix}0\\ \tau\end{bmatrix},
\label{eq:full_dynamics}
\end{equation}
第一行右端为零（基座不受控）；惯性块可由空间复合刚体算法以 \(O(n^{2})\) 复杂度计算。

设无外力矩且初始空间总动量为零；空间动量始终为零，经可逆的坐标变换后基座体坐标系动量也为零，故 \(H_b\dot x_b+H_{bm}\dot q=0\)，即
\begin{equation}
\dot x_b=-H_b^{-1}H_{bm}\dot q\triangleq J_{bm}(q)\dot q
\ \Longrightarrow\
\dot x_e=\bigl[J_m+J_bJ_{bm}\bigr]\dot q\triangleq J_g(q)\dot q ,
\label{eq:gjm}
\end{equation}
其前提是 \(H_b\) 在工作区域内一致正定。式~\eqref{eq:gjm}~表明基座运动是由关节运动诱导的反作用而非独立控制量。对 \(H_b\dot x_b+H_{bm}\dot q=0\) 求导并与 \eqref{eq:full_dynamics}~两行联立消去 \(\ddot x_b\)，得等效关节空间模型
\begin{equation}
M_e(g_b,q)\ddot q+c_e(g_b,q,\dot q)=\tau+d(t),\qquad
M_e(q)=H_m-H_{bm}^{\top}H_b^{-1}H_{bm},
\label{eq:ejm}
\end{equation}
其中 \(M_e\) 为 Schur 补，\(c_e\) 汇总递推动力学、动量约束求导与坐标变换产生的速度项，\(d(t)\) 为集总扰动。在选定坐标系下若 \(M_e\) 仅依赖关节、\(c_e=C_e(q,\dot q)\dot q\)，下文采用该简记；否则下列控制律中的 \(M_e,C_e\dot q\) 均须替换为完整的 \(M_e(g_b,q),c_e\)，并实时取得基座位姿以计算它们。将关节变量统一记为 \(q_m\)，后文使用的简记模型为
\begin{equation}
M_e(q_m)\ddot q_m+C_e(q_m,\dot q_m)\dot q_m=\tau+d(t).
\label{eq:model}
\end{equation}
式~\eqref{eq:model}~将原 \(6+n\) 维耦合系统降为 \(n\) 维。

\subsection{误差通道}
\label{sec:error_channel}

给定足够光滑的期望轨迹 \(q_d(t)\)，定义
\begin{equation}
e=q_m-q_d,\qquad \dot e=\dot q_m-\dot q_d,\qquad
z=\begin{bmatrix}e\\ \dot e\end{bmatrix}\in\R^{2n},
\label{eq:error_state}
\end{equation}
并取“计算力矩补偿 + 扰动补偿”结构
\begin{equation}
\tau=M_e(q_m)\bigl(\ddot q_d+\nu\bigr)+C_e(q_m,\dot q_m)\dot q_m-\hat d(t),
\label{eq:computed_torque}
\end{equation}
其中 \(\nu\in\R^{n}\) 为辅助加速度输入，\(\hat d\) 为扰动补偿量。

\begin{lemma}[误差通道解耦]\label{lem:linearization}
在假设~\ref{as:inertia}~下，将 \eqref{eq:computed_torque} 代入 \eqref{eq:model} 得
\begin{equation}
\ddot e_i(t)=\nu_i(t)+w_i(t),\qquad i=1,\dots,n,
\label{eq:decoupled_channel}
\end{equation}
其中 \(w=M_e^{-1}(q_m)(d-\hat d)\) 为残余加速度扰动。当 \(\hat d\equiv d\) 时 \(w\equiv0\)，误差通道化为 \(n\) 个解耦二阶积分器
\begin{equation}
\dot z=Az+B\nu,\quad
A=\begin{bmatrix}0&I_n\\ 0&0\end{bmatrix},\quad
B=\begin{bmatrix}0\\ I_n\end{bmatrix},
\label{eq:linear_error}
\end{equation}
且 \((A,B)\) 可控。
\end{lemma}

\begin{proof}
代入得 \(M_e\ddot q_m+C_e\dot q_m=M_e(\ddot q_d+\nu)+C_e\dot q_m-\hat d+d\)。由假设~\ref{as:inertia}~中 \(M_e\) 一致正定可左乘 \(M_e^{-1}\)，得 \(\ddot q_m=\ddot q_d+\nu+w\)，与 \eqref{eq:error_state} 相减即得 \eqref{eq:decoupled_channel}。可控性由 \(\operatorname{rank}[B,AB]=\operatorname{rank}\bigl[\begin{smallmatrix}0&I_n\\ I_n&0\end{smallmatrix}\bigr]=2n\) 保证。
\end{proof}

引理~\ref{lem:linearization}~是本结构的枢纽：它把强耦合的浮基跟踪问题化为逐关节标量二阶问题，使漏斗设计不需要反步、函数逼近与关节加速度测量，并使 PDF 理论可直接作用于 \eqref{eq:linear_error}。注意 \eqref{eq:decoupled_channel}~中的 \(w\) 并非待假设的上界：第~\ref{sec:ptdo}~节将证明 \(w\equiv\eta_2\)（PTDO 的第二误差分量），于是“\(t\ge T_o\) 后通道无扰动”是定理结论而非假设。

\subsection{性能函数与误差变换}
\label{sec:performance}

\begin{definition}[指定时间可达性能函数]\label{def:atpf}
称 \(\rho:\R_{\ge0}\to\R_{>0}\) 为可达时间 \(T_s\)、稳态半径 \(\rho_\infty\) 的指定时间可达性能函数，若 \textup{(i)} \(\rho\) 在 \([0,T_s)\) 上单调递减且为正；\textup{(ii)} \(\lim_{t\to T_s^-}\rho(t)=\rho_\infty>0\)；\textup{(iii)} \(\rho(t)\equiv\rho_\infty\) 于 \(t\ge T_s\)。
\end{definition}

\begin{definition}[误差变换函数]\label{def:transform}
称光滑严格递增的 \(S:\R\to(-1,1)\) 为误差变换函数，若 \(\lim_{\varepsilon\to\pm\infty}S(\varepsilon)=\pm1\)。对归一化误差 \(\xi_i=e_i/\rho_i\in(-1,1)\)，变换误差为 \(\varepsilon_i=S^{-1}(\xi_i)\)。
\end{definition}

对本文采用的 \(S=\tanh\)，\(S\) 是 \(\R\) 到 \((-1,1)\) 的微分同胚。对任意有限时间区间 \([0,T]\)，一致严格留在包络内部与变换误差一致有界等价（单个时刻的有限值不等于一致有界）：
\begin{equation}
\sup_{t\in[0,T]}\lvert e_i(t)/\rho_i(t)\rvert<1\iff\sup_{t\in[0,T]}\lvert\varepsilon_i(t)\rvert<\infty.
\label{eq:funnel_equiv}
\end{equation}
本文取 \(S(\varepsilon)=\tanh\varepsilon\)，但后文的证明直接在 \(\xi_i\) 坐标上进行（注~\ref{rem:barrier}~说明其必要性）。具体采用的多项式性能函数为
\begin{equation}
\rho_i(t)=
\begin{cases}
\rho_{i\infty}+\bigl(\rho_{i0}-\rho_{i\infty}\bigr)\varphi\!\Bigl(\dfrac{t}{T_o}\Bigr), & 0\le t<T_o,\\[6pt]
\rho_{i\infty}, & t\ge T_o,
\end{cases}
\qquad
\varphi(\tau)=1-\bigl(10\tau^{3}-15\tau^{4}+6\tau^{5}\bigr),
\label{eq:rho}
\end{equation}
其中 \(\varphi\) 为五次 smootherstep 的补函数，\(\rho_{i0}>\rho_{i\infty}>0\)。它满足定义~\ref{def:atpf}，并具有下述端点性质，这是把切换点取在 \(T_o\) 的原因。

\begin{lemma}[性能函数的端点性质]\label{lem:rho_smooth}
由 \eqref{eq:rho} 定义的 \(\rho_i\in C^{2}(\R_{\ge0})\)，且
\textup{(i)} \(\dot\rho_i(0)=0\)；\quad\textup{(ii)} \(\dot\rho_i(T_o)=\ddot\rho_i(T_o)=0\)；\quad
\textup{(iii)} \(\sup\lvert\dot\rho_i\rvert\le\frac{15}{8}\frac{\rho_{i0}-\rho_{i\infty}}{T_o}\)，\(\sup\lvert\ddot\rho_i\rvert\le\frac{10}{\sqrt3}\frac{\rho_{i0}-\rho_{i\infty}}{T_o^{2}}\)。
\end{lemma}

\begin{proof}
\(\varphi'(\tau)=-30\tau^{2}(1-\tau)^{2}\)、\(\varphi''(\tau)=-60\tau(1-\tau)(1-2\tau)\)，故 \(\varphi'(0)=\varphi'(1)=\varphi''(0)=\varphi''(1)=0\)；结合 \(\tau\to1^{-}\) 时 \(\rho_i\to\rho_{i\infty}\) 与右侧常值分支相接，得 (i)(ii) 与 \(C^{2}\)。再由 \(\max_{\tau}\lvert\varphi'\rvert=30/16=15/8\)（\(\tau=1/2\)）与 \(\max_{\tau}\lvert\varphi''\rvert=60/(6\sqrt3)=10/\sqrt3\)（\(\tau=(3\pm\sqrt3)/6\)）得 (iii)。
\end{proof}

\begin{remark}[为什么取 \(\dot\rho_i(0)=0\)]\label{rem:rho_zero_slope}
若改用 \(\rho_i\propto(1-t/T_o)^{m}\)，则 \(\lvert\dot\rho_i\rvert\) 在 \(t=0\) 处\emph{取最大值} \(m(\rho_{i0}-\rho_{i\infty})/T_o\)。此时初始速度误差为零而包络立即收缩，由 \eqref{eq:xi_s} 得 \(\dot\xi_i(0)=-\xi_i(0)\dot\rho_i(0)/\rho_{i0}\) 与 \(\xi_i(0)\) 同量级，控制律 \eqref{eq:nu_ppc} 于是在 \(t=0\) 产生与 \(\rho_{i0}\) 同阶的加速度指令；在等效惯量达 \(3\times10^{2}\,\mathrm{kg\cdot m^2}\) 的关节上，这直接触发饱和（本文数值实验中初始力矩即达 \(20\,\mathrm{N\cdot m}\) 限幅）。取 \(\varphi'(0)=0\) 使包络“从静止开始收缩”，把初始力矩冲击降为 \(O(k_1k_2\xi_i\rho_{i0})\)；而 \(\varphi'(1)=\varphi''(1)=0\) 又保证切换点不产生由 \(\rho\) 引起的力矩跳变（注~\ref{rem:jump}）。
\end{remark}

\subsection{控制问题与假设}
\label{sec:problem}

\begin{problem}[双重预设跟踪]\label{prob:dual}
给定性能函数 \eqref{eq:rho}、预设观测时间 \(T_o\) 与人工延迟 \(h\)，设计形如 \eqref{eq:computed_torque} 的力矩律使
\textup{(i)} \textbf{包络} \(\lvert e_i(t)\rvert<\rho_i(t)\)，\(\forall t\ge0\)；
\textup{(ii)} \textbf{进带} \(\lvert e_i(T_o)\rvert\le\rho_{i\infty}\)（预设时间 \(T_o\) 进入预设误差带）；
\textup{(iii)} \textbf{归零} \(e(t)=\dot e(t)=0\)，\(\forall t\ge T_o+2h\)（再经预设时间 \(2h\) 精确归零）。
理想条件下要求 (i)--(iii) 同时成立；当周期零化与补偿存在有界误差时，只在满足命题~\ref{prop:approx} 的递推条件下给出含初态瞬态项的有限时间界及最终界，不能把最终界写成全时域界。
\end{problem}

问题~\ref{prob:dual}~(ii)(iii) 分别由 PPC 与 PDF 承担，二者共享交接时刻 \(T_o\)，总预设时间 \(T_{\mathrm{pre}}=T_o+2h\)。

\begin{assumption}[模型与惯性]\label{as:inertia}
\(M_e\) 对称、连续、一致正定（\(\exists\underline m>0,\ M_e\succeq\underline m I_n\)）；实际控制所用的 \(M_e\) 与速度项 \(C_e\dot q_m\) 精确已知；若其依赖基座位姿，则基座位姿可测或可重构。
\end{assumption}

\begin{assumption}[可测性与参考轨迹]\label{as:meas}
\(q_m(t)\)、\(\dot q_m(t)\) 可实时测量；\(q_d,\dot q_d,\ddot q_d\) 连续有界。关节加速度 \(\ddot q_m\) 与 \(d(t)\) 的数值\emph{不}需要测量或已知。
\end{assumption}

\begin{assumption}[扰动有界性]\label{as:d}
\(\Delta_a=M_e^{-1}(q_m)d(t)\) 及其导数 \(\dot\Delta_a\) 一致有界，但上界数值未必已知。
\end{assumption}

\begin{assumption}[初始条件]\label{as:ic}
\(\lvert e_i(0)\rvert<\rho_{i0}\)，\(i=1,\dots,n\)。
\end{assumption}

\begin{assumption}[精确归零]\label{as:ideal}
阶段 II 的精确结论要求同时满足：\textup{(i)} 连续时间实现；\textup{(ii)} 无执行器饱和；\textup{(iii)} 模型补偿精确；\textup{(iv)} PTDO 存在注~\ref{rem:certificate_gap}~中的收敛证书 \((\star)\)。前三条不满足时结论退化为命题~\ref{prop:approx}~的显式有界性。
\end{assumption}

\begin{assumption}[观测器条件]\label{as:observer}
对所选 \(T_c\le T_o\)、\(\eta\in(0,1)\)、\(\sigma>0\)，需分别核验两个层次：
\textup{(A)} 误差系统 \eqref{eq:observer_errordyn} 在 \([0,T_o]\) 上存在一致有界的 Filippov 解（等价于无有限时刻逃逸，可由不变集或仿真监控核验）；
\textup{(B)} 增益条件 \(\lvert\dot\Delta_{a,i}\rvert\le\kappa_i<\pi/(\eta T_c)\)。
这两条足以支撑阶段 I 的包络与进带结论（定理~\ref{thm:ptdo_core}、定理~\ref{thm:algo}~(ii)(iii)），且不需任何额外假设。若还要阶段 II 的\emph{精确}补偿，则需另提供针对 \eqref{eq:ptdo}--\eqref{eq:phi} 的收敛证书 \((\star)\)；本文不提供该证书，其缺失原因与可选修补见注~\ref{rem:certificate_gap}。
\end{assumption}

\begin{remark}\label{rem:assumption_cmp}
与“扰动及其导数有界且\emph{上界已知}”的常规设计相比，假设~\ref{as:d}~只要求有界。后文将证明该弱化不削弱包络结论（定理~\ref{thm:funnel}），只削弱进带\emph{深度}的定量性（注~\ref{rem:unknown_w}）。\end{remark}
